% ============================================================================
% Fichier  : partie4/P04_C24_audit_tiers_nuage.tex
% Titre    : Auditer les Tiers et le Nuage
% Auteur   : MINKA MI NGUIDJOÏ Thierry Emmanuel
% Version  : 1.0 -- novembre 2026
% ============================================================================
% Base : ACSI C17 (section prestataires et nuage ; contrat DIGICORE, huit
% clauses 4/3/1) et ACSI C10 (rapport ISAE 3402 de type 1 de DIGICORE :
% 12 contrôles cotés NEK, µ 0, U 1, H 0).
% Sources : IAASB ISAE 3402 ; NIST SP 800-145 (définition de
% l'informatique en nuage) ; ISO/IEC 27017:2015 et 27018:2019 (mesures pour
% le nuage et les données personnelles dans le nuage) ; Cloud Security
% Alliance, Cloud Controls Matrix v4 ; loi n° 2024/017 (données
% personnelles, à vérifier au Journal officiel) ; COBAC R-2016/04.
% Faits nouveaux (LiZbiZ) : prestataire d'exploitation de Garoua et
% Bafoussam depuis septembre 2026 (cause de la hausse de U du domaine E).
% Corrigés : corriges/Corr_C24.tex.
% ============================================================================

\chapter{Auditer les Tiers et le Nuage}
\label{chap:audit-tiers}
\chaptermark{Tiers et nuage}

\epigraph{%
    « On externalise une activité.\\
    On n'externalise jamais sa responsabilité. »%
}{Principe de l'externalisation}

\niveauChapitre{Intermédiaire}{%
    Maîtrise des Chapitres~\ref{chap:normes} (ISAE~3402),
    \ref{chap:audit-smsi} (SMSI) et~\ref{chap:audit-continuite}
    (continuité)%
}

\begin{boxobjectifs}
\begin{itemize}
    \item Identifier les prestataires critiques et évaluer la dépendance.
    \item Auditer un contrat de prestation clause par clause.
    \item Utiliser un rapport ISAE~3402 de type~1 ou~2, et ses contrôles
          complémentaires.
    \item Définir les modèles de service en nuage et répartir les
          responsabilités entre fournisseur et client.
    \item Vérifier la localisation des données et la réversibilité.
    \item Tenir un registre des prestataires critiques.
\end{itemize}
\end{boxobjectifs}

\bigskip

% ============================================================================
\section*{Question directrice : comment auditer ce qu'on ne voit pas~?}
% ============================================================================

Le système central de la BANQUE SANAGA est exploité par DIGICORE AFRIQUE.
Sa messagerie est un service en nuage depuis 2025. Chez LiZbiZ, les centres
de Garoua et de Bafoussam sont exploités depuis septembre~2026 par un
prestataire local, et l'incertitude du domaine Exploitation a augmenté
d'autant (Chapitre~\ref{chap:suivi}). Dans les trois cas, l'auditeur ne peut
pas entrer dans les locaux du prestataire, et ne lit pas ses journaux.
Comment conclure~?

\medskip
\textbf{Plan du chapitre.}
\begin{enumerate}
    \item Les prestataires critiques et la dépendance
          (section~\ref{sec:c24-critiques}).
    \item Auditer un contrat (section~\ref{sec:c24-contrat}).
    \item Les rapports d'assurance (section~\ref{sec:c24-isae}).
    \item Le nuage et la responsabilité partagée (section~\ref{sec:c24-nuage}).
    \item La localisation des données et la réversibilité
          (section~\ref{sec:c24-donnees}).
\end{enumerate}

\separateur

% ============================================================================
\section{Les Prestataires Critiques}
\label{sec:c24-critiques}
% ============================================================================

\subsection{Identifier}

\begin{boxdef}
\termecle{Prestataire critique}\index{prestataire critique}~: prestataire
dont la défaillance interromprait un service critique de l'organisation,
ou exposerait des données sensibles, au-delà de l'appétence au risque.
\end{boxdef}

Pour une banque de la CEMAC, l'externalisation d'une activité essentielle
reste couverte par le dispositif de contrôle interne exigé par la
COBAC~\cite{cobac2016r04}~: le risque opérationnel d'un prestataire est un
risque de la banque.

\subsection{Évaluer la dépendance}

\begin{tableaucours}{L{3.2cm}L{4.8cm}L{4.2cm}}{%
    \tableauHeader{Question} &
    \tableauHeader{Ce qu'elle révèle} &
    \tableauHeader{Exemple (SANAGA)}
}
Quel service en dépend~? & L'impact d'une défaillance & Système central~:
toute l'activité \\
Existe-t-il une alternative~? & La possibilité de changer & Aucune
alternative connue \\
En combien de temps~? & La durée de la dépendance & Plus de 12~mois pour
un nouveau prestataire \\
Le prestataire sous-traite-t-il~? & La chaîne réelle des acteurs & Inconnu
(clause absente) \\
\end{tableaucours}
\vspace{-0.8em}
\begin{center}{\small\itshape Tableau~24.1. Évaluer la dépendance à un
prestataire}\end{center}

% ============================================================================
\section{Auditer un Contrat}
\label{sec:c24-contrat}
% ============================================================================

\subsection{Les clauses}

\begin{tableaucours}{L{4.0cm}L{8.2cm}}{%
    \tableauHeader{Clause} &
    \tableauHeader{Ce que l'auditeur vérifie}
}
Niveaux de service & Disponibilité, délais, RTO et RPO alignés sur ceux de
l'organisation~; mesure définie~; pénalités appliquées \\
Droit d'audit & L'organisation, son auditeur et son régulateur peuvent
auditer le prestataire, ou obtiennent un rapport d'assurance équivalent \\
Sécurité et incidents & Mesures exigées~; notification des incidents dans
un délai fixé \\
Localisation des données & Lieu de stockage connu, compatible avec la loi
sur les données personnelles (Chapitre~\ref{chap:droit}) \\
Sous-traitance ultérieure & Recours à d'autres prestataires soumis à
accord \\
Réversibilité & Restitution des données dans un format exploitable, et
assistance à la sortie \\
Continuité du prestataire & Plan de continuité testé, dont les résultats
sont communiqués \\
Confidentialité & Engagement du prestataire et de son personnel \\
\end{tableaucours}
\vspace{-0.8em}
\begin{center}{\small\itshape Tableau~24.2. Clauses à auditer dans un
contrat de prestation ou de service en nuage}\end{center}

\begin{boxscenario}{Opération SANAGA. Le contrat de DIGICORE AFRIQUE}
\textbf{Faits.} Mme~EKAMBI lit le contrat de maintenance du système central
conclu avec DIGICORE AFRIQUE. Il prévoit des niveaux de service et des
pénalités, une clause de confidentialité, la localisation des données à
Douala et un plan de continuité du prestataire. Il ne prévoit ni droit
d'audit, ni réversibilité, ni encadrement de la sous-traitance ultérieure.
L'annexe sur la notification des incidents, citée par le contrat, n'a pas
été remise.

\medskip
\textbf{Cotation} des huit clauses~: 4~\OUI{}, 3~\NON{}, 1~\NEK{}.

\medskip
\textbf{Analyse} (confiance élevée). Sans droit d'audit, la banque dépend du
rapport que le prestataire choisit de produire~: un type~1
(section~\ref{sec:c24-isae}). Sans réversibilité, elle ne peut pas changer de
prestataire sans risque. La recommandation porte sur un avenant, à
négocier avant le renouvellement.
\end{boxscenario}

\begin{boxalert}
\textbf{Violation typique.} Coter \OUI{} la clause de niveaux de service
parce qu'elle existe, sans vérifier que les pénalités ont été appliquées
lors de la panne.

\smallskip
\textbf{Règle~:} une clause se juge sur sa conception (existe-t-elle,
est-elle mesurable~?) puis sur son exécution (a-t-elle été appliquée quand
elle devait l'être~?). Après 72,5~heures de panne, l'absence de pénalité
réclamée est un constat.
\end{boxalert}

% ============================================================================
\section{Les Rapports d'Assurance}
\label{sec:c24-isae}
% ============================================================================

\subsection{Type 1 et type 2}

Un prestataire qui sert de nombreux clients ne peut pas accueillir
l'auditeur de chacun. Il fait réaliser, par un auditeur indépendant, un
rapport sur ses propres contrôles, selon la norme ISAE~3402~\cite{iaasb2011isae3402}
(Chapitre~\ref{chap:normes}).

\begin{tableaucours}{L{2.4cm}L{5.0cm}L{4.8cm}}{%
    \tableauHeader{Type} &
    \tableauHeader{Ce qu'il atteste} &
    \tableauHeader{Usage pour l'auditeur du client}
}
Type~1 & La description des contrôles et leur conception, à une date &
Comprendre le dispositif~; ne permet pas de coter l'efficacité \\
Type~2 & Conception et efficacité des contrôles sur une période (souvent
6 à 12 mois) & Coter l'efficacité sur la période couverte, sous réserve des
écarts signalés \\
\end{tableaucours}
\vspace{-0.8em}
\begin{center}{\small\itshape Tableau~24.3. Rapports ISAE~3402 de type~1 et
de type~2}\end{center}

\begin{boxlecture}{Le rapport de type 1 de DIGICORE AFRIQUE}
Le rapport décrit douze contrôles et atteste leur conception à la date du
31~décembre~2025. Il ne teste aucune exécution.

\medskip
\textbf{Constatation.} Pour la question « le contrôle a-t-il fonctionné en
2026~? », les douze contrôles sont \NEK{}~: $\mu = 0$~; $U = 1$~; $H = 0$.

\medskip
\textbf{Analyse} (confiance élevée). C'est le cas limite de l'ignorance
unanime (Chapitre~\ref{chap:conclure})~: l'entropie nulle ne traduit
aucune certitude, mais l'absence complète d'information sur l'efficacité.
La banque doit obtenir un rapport de type~2, ou un droit d'audit.

\medskip
\hfill\lecture{0{,}00}{1{,}00}{0{,}00}
\end{boxlecture}

\subsection{Les contrôles complémentaires}

Tout rapport ISAE~3402 suppose que le client exerce certains contrôles de
son côté~: par exemple, revoir chaque trimestre la liste des comptes que
le prestataire a ouverts à ses agents. Ces \termecle{contrôles
complémentaires}\index{contrôles complémentaires} restent à tester chez le
client~; sans eux, le rapport le plus favorable ne couvre pas le risque.

% ============================================================================
\section{Le Nuage et la Responsabilité Partagée}
\label{sec:c24-nuage}
% ============================================================================

\subsection{Trois modèles de service}

Le NIST définit l'informatique en nuage par cinq caractéristiques
(service à la demande, accès par le réseau, mutualisation des ressources,
élasticité, service mesuré) et trois modèles de service~\cite{nist2011sp800145}.

\begin{tableaucours}{L{2.4cm}L{4.4cm}L{5.4cm}}{%
    \tableauHeader{Modèle} &
    \tableauHeader{Le fournisseur gère} &
    \tableauHeader{Le client reste responsable de}
}
IaaS (infrastructure) & Matériel, réseau, virtualisation & Systèmes,
applications, données, comptes, configuration de sécurité \\
PaaS (plateforme) & En plus~: systèmes et intergiciels & Applications,
données, comptes, paramètres \\
SaaS (logiciel) & En plus~: l'application & Données, comptes et droits,
paramètres de sécurité de l'application \\
\end{tableaucours}
\vspace{-0.8em}
\begin{center}{\small\itshape Tableau~24.4. Les modèles de service et la
responsabilité partagée}\end{center}

Quel que soit le modèle, le client reste responsable de ses données, de
ses comptes et de leurs droits. L'auditeur établit, pour chaque contrôle,
qui en est responsable~; un contrôle que chacun croit à la charge de
l'autre n'est assuré par personne.

\subsection{Les référentiels du nuage}

L'ISO/IEC~27017 complète l'ISO/IEC~27002 par des mesures propres au
nuage~\cite{iso27017_2015}~; l'ISO/IEC~27018 traite de la protection des
données personnelles chez un fournisseur de nuage public~\cite{iso27018_2019}.
La Cloud Security Alliance publie une matrice de contrôles, la \emph{Cloud
Controls Matrix}, qui précise pour chaque contrôle la répartition entre
fournisseur et client~\cite{csa2021ccm}.

\begin{boxscenario}{Opération SANAGA. La messagerie dans le nuage}
\textbf{Faits.} La messagerie de la banque est un service en nuage (SaaS)
depuis 2025. Le fournisseur produit un rapport de type~2 sur son
infrastructure. La banque n'a jamais activé la double authentification
pour ses comptes, ni configuré la conservation des journaux d'accès.

\medskip
\textbf{Analyse} (confiance élevée). Le rapport du fournisseur est
favorable, mais il ne couvre pas ce qui relève de la banque~: la gestion
des comptes et la configuration. Les deux contrôles sont \NON{}, et ils
sont de la responsabilité de la banque.
\end{boxscenario}

% ============================================================================
\section{Localisation des Données et Réversibilité}
\label{sec:c24-donnees}
% ============================================================================

\subsection{Où sont les données~?}

La loi n°~2024/017 encadre le traitement et le transfert des données
personnelles~; les conditions d'un transfert hors du Cameroun sont à
vérifier dans le texte publié au Journal officiel. L'auditeur vérifie que
le contrat indique le lieu de stockage et de traitement, que ce lieu est
compatible avec la loi, et que le prestataire ne peut pas le changer sans
accord.

\subsection{Peut-on partir~?}

\begin{boxprocedure}
\textbf{Auditer la réversibilité}
\begin{enumerate}
    \item Le contrat prévoit-il la restitution de toutes les données, dans
          un format documenté et exploitable, dans un délai fixé~?
    \item Prévoit-il une assistance à la reprise par un autre prestataire,
          et son coût~?
    \item La destruction des données chez le prestataire, après
          restitution, est-elle prévue et attestée~?
    \item Un plan de réversibilité existe-t-il, et a-t-il été testé, au
          moins par une restitution partielle~?
\end{enumerate}
\end{boxprocedure}

\subsection{Le registre des prestataires critiques}

\begin{boxmetier}{Gestionnaire de risques}
\textbf{Tenir le registre des prestataires critiques}
\begin{enumerate}
    \item Recenser les prestataires dont dépend un service critique, et
          le service concerné.
    \item Pour chacun, noter les clauses présentes, le dernier rapport
          d'assurance reçu, son type et sa période.
    \item Évaluer la dépendance~: existe-t-il une solution de remplacement,
          en combien de temps~?
    \item Porter au registre des risques tout prestataire critique sans
          droit d'audit ni réversibilité.
    \item Revoir le registre chaque année, et avant chaque renouvellement
          de contrat.
\end{enumerate}
\end{boxmetier}

\begin{boxmetier}{Responsable de la conformité}
\textbf{Vérifier les transferts de données.} Le responsable de la
conformité tient la liste des traitements confiés à des prestataires, avec
le lieu de stockage et le fondement du transfert. Un prestataire en nuage
qui ne peut pas dire où sont les données est un risque de conformité, quel
que soit son niveau de sécurité.
\end{boxmetier}

\begin{boxscenario}{Opération LIZBIZ. Le prestataire de Garoua}
\textbf{Faits.} Depuis septembre~2026, un prestataire exploite les centres
de Garoua et de Bafoussam. Le contrat, de trois pages, ne prévoit ni
rapport d'activité, ni droit d'audit, ni notification des incidents.

\medskip
\textbf{Analyse} (confiance élevée). La hausse de l'incertitude du domaine
Exploitation (de 0,30 à 0,40, Chapitre~\ref{chap:suivi}) a ici sa cause~:
LiZbiZ a externalisé l'exploitation sans se donner les moyens d'en savoir
quelque chose. La recommandation~: un avenant (rapport mensuel, droit
d'audit, notification des incidents sous 24~heures), et un interlocuteur
interne chargé de suivre le prestataire.
\end{boxscenario}

\begin{boxinfo}
\textbf{Ce que dit la recherche.} Le premier axiome d'A0, l'observabilité
partielle~\cite{minka2026auditologie}, est ici à son extrême~: l'auditeur
du client n'observe le prestataire qu'à travers ce que le contrat lui donne
le droit de voir. Le taux $U$ mesure directement la qualité de ce droit~:
un contrat sans droit d'audit produit mécaniquement des \NEK{}.
\end{boxinfo}

\separateur

\begin{boxresume}
\textbf{Ce qu'il faut retenir de ce chapitre~:}
\begin{itemize}
    \item \textbf{On externalise une activité, pas sa responsabilité}~; un
          prestataire critique se recense et sa dépendance s'évalue.
    \item \textbf{Un contrat s'audite clause par clause}, en conception
          puis en exécution.
    \item \textbf{Un rapport ISAE~3402 de type~1} ne dit rien de
          l'efficacité~; le type~2 couvre une période~; les contrôles
          complémentaires restent à tester chez le client.
    \item \textbf{Dans le nuage}, la responsabilité est partagée selon le
          modèle (IaaS, PaaS, SaaS)~; le client reste responsable de ses
          données, de ses comptes et de sa configuration.
    \item \textbf{La localisation des données et la réversibilité} se
          vérifient dans le contrat et, pour la réversibilité, par un test.
\end{itemize}
\end{boxresume}

\begin{boxinfo}
\textbf{Points de contrôle de l'élève.} Je sais~: identifier un prestataire
critique~; auditer les huit clauses d'un contrat~; dire ce que vaut un
rapport de type~1 ou~2~; répartir les responsabilités selon le modèle de
nuage~; auditer la réversibilité~; tenir un registre des prestataires.
\end{boxinfo}

\bigskip

\begin{boxexo}[title={\ding{45}\quad Exercice~24.1. Critique ou non~? \niveaubase}]
Classez comme critiques ou non, en justifiant~: (a) le fournisseur de
fournitures de bureau de LiZbiZ~; (b) l'hébergeur du portail de suivi des
expéditions~; (c) le prestataire de nettoyage du centre de données~; (d)
l'éditeur de l'application intégrée.
\end{boxexo}

\begin{boxexo}[title={\ding{45}\quad Exercice~24.2. Coter un contrat \niveaumoyen}]
Cotez les huit clauses du Tableau~24.2 pour le contrat du prestataire de
Garoua (scénario de la section~\ref{sec:c24-donnees}), calculez $\mu$,
$U$ et $H$, et rédigez la recommandation principale.
\end{boxexo}

\begin{boxexo}[title={\ding{45}\quad Exercice~24.3. Type 1 ou type 2 \niveaubase}]
Un rapport ISAE~3402 de type~2 couvre la période du 1\textsuperscript{er}
avril au 30~septembre~2026 et signale un écart sur la gestion des
changements. Que pouvez-vous coter pour 2026, et que manque-t-il~?
\end{boxexo}

\begin{boxexo}[title={\ding{45}\quad Exercice~24.4. Les contrôles complémentaires \niveaumoyen}]
Le rapport de type~2 du fournisseur de messagerie liste trois contrôles
complémentaires attendus du client~: revue des comptes, activation de la
double authentification, conservation des journaux. Comment les testez-vous
à la banque, et quelle cote leur donnez-vous d'après le scénario de la
section~\ref{sec:c24-nuage}~?
\end{boxexo}

\begin{boxexo}[title={\ding{45}\quad Exercice~24.5. Qui est responsable~? \niveaumoyen}]
Pour un service PaaS qui héberge l'application de paie de l'ANPP, dites
qui est responsable de~: les correctifs du système d'exploitation~; les
droits des agents de l'ANPP~; le chiffrement des données au repos (selon le
contrat)~; la sauvegarde de la base~; le code de l'application.
\end{boxexo}

\begin{boxexo}[title={\ding{45}\quad Exercice~24.6. L'ignorance unanime \niveaumoyen}]
Expliquez au comité d'audit, en cinq lignes, pourquoi la lecture du
rapport de type~1 de DIGICORE ($H = 0$) ne signifie pas que « tout est
clair ».
\end{boxexo}

\begin{boxexo}[title={\ding{45}\quad Exercice~24.7. Les pénalités \niveaumoyen}]
Le contrat de DIGICORE prévoit une pénalité de 2~\% de la redevance
mensuelle par heure d'indisponibilité au-delà de 4~heures, plafonnée à
30~\%. Calculez la pénalité due pour la panne de mars~2026, et dites ce que
signifie, pour l'auditeur, le fait qu'elle n'ait pas été réclamée.
\end{boxexo}

\begin{boxexo}[title={\ding{45}\quad Exercice~24.8. La réversibilité \niveaumoyen}]
Rédigez la clause de réversibilité que vous recommandez pour l'avenant au
contrat de DIGICORE, et le test qui permettrait de vérifier qu'elle est
applicable.
\end{boxexo}

\begin{boxexo}[title={\ding{45}\quad Exercice~24.9. Un nuage pour l'ANPP \niveauexpert}]
L'ANPP envisage d'héberger le portail de consultation des bulletins de paie
chez un fournisseur de service en nuage.
\begin{enumerate}[(a)]
    \item Dressez la liste des contrôles et, pour chacun, dites s'il
          relève du fournisseur ou de l'ANPP.
    \item Quelles clauses le contrat doit-il contenir~? Lesquelles sont
          liées au droit camerounais~?
    \item Le fournisseur propose un rapport ISAE~3402 de type~1. Que
          demandez-vous, et pourquoi~?
    \item Comment la Chambre des comptes pourra-t-elle auditer ce service~?
\end{enumerate}
\end{boxexo}

\begin{boxexo}[title={\ding{45}\quad Exercice~24.10. Le registre de la BANQUE SANAGA \niveauexpert}]
Construisez le registre des prestataires critiques de la BANQUE SANAGA pour
trois prestataires (DIGICORE, messagerie, monétique), avec les colonnes de
la procédure, et proposez les deux risques à porter au registre des
risques.
\end{boxexo}

\bigskip

\begin{boxinfo}
\textbf{Pour aller plus loin.}
\begin{itemize}
    \item IAASB, ISAE~3402~\cite{iaasb2011isae3402}.
    \item NIST, SP~800-145~\cite{nist2011sp800145}~: la définition du nuage.
    \item ISO/IEC~27017~\cite{iso27017_2015} et
          ISO/IEC~27018~\cite{iso27018_2019}.
    \item Cloud Security Alliance, \textit{Cloud Controls Matrix}~\cite{csa2021ccm}.
\end{itemize}
\end{boxinfo}

\begin{boxinfo}
\textbf{Transition.} Les tiers posent souvent une question de conformité~:
où sont les données, et la loi est-elle respectée~? Le
Chapitre~\ref{chap:audit-conformite} présente l'audit de conformité,
méthode complète, du registre des obligations au rapport.
\end{boxinfo}
